This document must address the following key questions:
\begin{itemize}
    \item What is the specific subject of study?
    \item Why is this subject important? 
    \item Who is the intended audience?
    \item What impact should the project have on their understanding?
    \item How will the project achieve its communication objectives?
    \item Descrizione dettagliata dei diversi punti da affrontare
    \item Quali sono i risultati attesi?
    \item Quali metodologie utilizzare per costruire dati sintetici?
    \item ?Analisi rappresentatività? (perdita di informazione e problematiche embedding (Probing MDL????))
    \item Scelta di un dominio
    \item Analisi sui dati reali e risutlati attesi (Bayesian statisrtics)
    \item Implicazioni in termini di sicurezza sul dominio
\end{itemize}

This investigation formalizes the phenomenon of \textit{structural aliasing} within patch-based time-series foundation models, focusing specifically on Amazon's Chronos-Bolt framework.

\section{Purpose}
This project deploys a highly localized, layer-resolved evaluation blueprint across five distinct stages of the transformer architecture, probing immediately after individual encoder and decoder blocks, and directly following the final direct quantile output projection layer.

By mapping representations across these specific layers, the project isolates where task-relevant signal attributes are preserved or compromised under architectural synchronization constraints. The investigation evaluates the model using a custom hyperparameter configuration, where patch size is held constant at $P=16$, and the stride is systematically modified within the range of $S \in [8, 16]$ to explore the explicit impact of patch geometry on aliasing effects.

Ultimately, this study aims to provide a comprehensive understanding of the relationships between raw signal characteristics and the latent representations learned by the model, systematically identifying the critical links between sampling frequency ($f_s$), patch size ($P$), and stride ($S$) in order to propose actionable, principled architectural solutions to mitigate the phenomenon.

\section{Intended Audience and Target Impact}
This analysis is directly intended for researchers and practitioners utilizing large language models within the Chronos family for time-series analysis, representation learning, and forecasting, as well as individuals interested in characterizing the fundamental limitations, structural boundaries, and latent capabilities of these architectures. The core objective is to challenge the baseline assumption that pre-trained temporal foundation models inherently preserve frequency signal attributes when processing continuous data tracks, provided they satisfy classical sampling theorems. 